\section{混合\texorpdfstring{$A_2$}{A2}によるheartのずれの具体化}\label{sec:A2}
この節は定理\ref{thm:derived}のgenuine楕円gluingで，$q$非torsionとする．
$H=-K$はnef，$H^2>0$．楕円境界と交わらない一次元以下の対象を
$\cZ=\{M\in\AR:\dim M\le1,\ F_qM=0\}$と書く．

\subsection{単純根sheafと有限長部分圏}
\begin{lemma}\label{lem:offboundary}
$\cZ$は有限長Serre部分圏であり，その既約純一次元root sheafは$\AR$の単純対象である．
\end{lemma}
\begin{proof}
境界非交叉な対象の零次元商$Z$には$i^*Z=0$が従い，$\theta Z\to Z$は全射となる．
両者の数値類は等しいので\cite[Corollary 7.7]{Rains}から同型である．
$[Z]=\ell[\mathrm{pt}]$，$\ell\ge0$に対して境界determinantは$q^{\ell}$であり，
非torsion性から$\ell=0$，従って$Z=0$．既約純一次元sheafの非零部分対象の商は
零次元なので，この場合には単純性が得られる．

$0\to N\to M\to Q\to0$，$M\in\cZ$とすると，nef性と加法性から
$Hc_1(N)=Hc_1(Q)=0$．$i^*Q=0$なので$\theta Q\to Q$は全射であり，
rank零，$Kc_1(Q)=0$と\cite[Corollary 7.15]{Rains}から数値類は変わらない．
Corollary 7.7と自然性により$\theta Q\to Q$，$\theta N\to N$はいずれも同型で，
$N,Q\in\cZ$となる．拡大閉性は境界triangleから従う．
真の降鎖の逐次商は零次元になれないため，\cite[Proposition 7.32]{Rains}の
finitely partitive性が降鎖を止める．coherent heartのnoetherian性と合わせて有限長となる．
\end{proof}

\begin{corollary}[全境界外圏のtorsion pair]\label{cor:Ztilt}
定理\ref{thm:globaltilt}のtorsion pairは$\cZ$へ制限され，
\[
\cT\cap\cZ=\ker\Ext^1(e,-),\qquad
\cF\cap\cZ=\ker\Hom(e,-)
\]
となる．これは二根部分圏だけでなく全$\cZ$についての結論である．
\end{corollary}
\begin{proof}
$\cZ$がSerre部分圏なので，その対象のtorsion分解の部分対象と商も$\cZ$に属する．
$M\in\cZ$では$g$が可逆に作用し，各$\Ext^j(e,s^nM)$の移行写像が同型となる．
その局所化を$V^j(M)$とすれば$\dim V^j(M)=\dim\Ext^j(e,M)$であり，
$H^j\Phi(M)=0$と$\Ext^j(e,M)=0$は同値である．
\end{proof}

異なる二根sheaf$S_\alpha,S_\beta\in\cZ$には
\begin{equation}\label{eq:rootext}
\Hom(S_\alpha,S_\beta)=\Ext^2(S_\alpha,S_\beta)=0,
\qquad \dim\Ext^1(S_\alpha,S_\beta)=\alpha\cdot\beta
\end{equation}
が成り立つ．実際，単純性からHomが消え，境界非交叉とSerre双対性からExt$^2$も消える．
rank-zero Euler形式$\chi(S_\alpha,S_\beta)=-\alpha\beta$が残る次元を与える．

隣接根$\alpha\beta=1$，$\chi(e,S_+)=p>0$，$\chi(e,S_-)=-n<0$を固定し，
両向きの非分裂拡大を
\begin{equation}\label{eq:PQ}
0\to S_-\to P\to S_+\to0,
\qquad 0\to S_+\to Q\to S_-\to0
\end{equation}
とする．存在は\eqref{eq:rootext}から従い，各中間項は同型を除き一意である．

\begin{proposition}[connecting mapの最大rank]\label{prop:rank}
$P$のconnecting map$\delta:\Hom(e,S_+)\to\Ext^1(e,S_-)$は
$\rk\delta=\min(p,n)$を満たす．従って
\[
(h^0(P),h^1(P))=((p-n)_+,(n-p)_+),\qquad (h^0(Q),h^1(Q))=(p,n),
\]
ただし$h^j(M)=\dim\Ext^j(e,M)$，$x_+=\max(x,0)$とする．
このrankの証明には定理\ref{thm:derived}を必要としない．
\end{proposition}
\begin{proof}
符号別Extから$\Hom(e,P)=\ker\delta$，$\Ext^1(e,P)=\operatorname{coker}\delta$．
非零$f:e\to P$があれば，$P\to S_+$との合成は非零で全射となる．
像と$S_-$との交わりは零か$S_-$であり，前者なら拡大が分裂するので後者である．
従って$f$は全射．核$L$はrank-one torsion-freeで$c_1(L)=-(\alpha+\beta)$．
Corollary 10.12と$H(K+\alpha+\beta)=-H^2<0$により
$\Ext^2(e,L)=\Hom(L,\theta e)^*=0$．$e$のexceptionalityから$h^1(P)=0$となる．
$p>n$なら$\ker\delta\ne0$なので$\delta$は全射．$p\le n$で$h^0(P)>0$とすると
$\chi(e,P)=h^0(P)>0$に矛盾するので，この場合は単射である．
$Q$には逆向き短完全列の六項完全列を使う．
\end{proof}

\begin{proposition}[二根の全有限長対象]\label{prop:A2category}
$\cC=\langle S_+,S_-\rangle_{\mathrm{Serre}}$は
$\operatorname{mod}_{\mathrm{fd}}\Pi(A_2)$と同値であり，直既約対象は$S_+,S_-,P,Q$の四つである．
ここで$\Pi(A_2)$は二頂点・両向き二矢と，両方の長さ二の道を零とする関係を持つ代数である．
$p>n$では$(\cT\cap\cC,\cF\cap\cC)=(\add(S_+,P),\add(S_-))$，
$p<n$では$(\add(S_+),\add(S_-,P))$である．
\end{proposition}
\begin{proof}
非分裂性と単純性から，順序$(S_+,S_-,P,Q)$でHom次元は
\[
\begin{pmatrix}
1&0&0&1\\0&1&1&0\\1&0&1&1\\0&1&1&1
\end{pmatrix}.
\]
例えば$\Hom(P,Q)$は$P\twoheadrightarrow S_+\hookrightarrow Q$で生成され，
逆方向も同様で，両合成は零である．$\theta$はこれらの対象で自然に恒等なので
\[
\dim\Ext^1(U,V)=\dim\Hom(U,V)+\dim\Hom(V,U)+c_1(U)c_1(V).
\]
上の表を代入すると，$P,Q$から二単純へのExt$^1$はすべて零である．
長さによる帰納法から$P\oplus Q$は$\cC$のprojective generatorであり，
そのEnd代数は表示した四次元代数となる．
その表現$x:U_+\to U_-$，$y:U_-\to U_+$，$xy=yx=0$は，
$\operatorname{im}y\subset\ker x$，$\operatorname{im}x\subset\ker y$に適合する基底により，
一方の矢だけが非零の一次元対と孤立一次元空間へ分解される．従って直既約は四つである．
最後の主張は命題\ref{prop:rank}と符号別Extによる．
\end{proof}

$\cC$の内在的ExtはambientのExtと異なる．実際$\Pi(A_2)$の単純加群には
二周期的なprojective resolutionがあり，$\Ext^{2j}_{\cC}(S_+,S_+)=k$ $(j\ge0)$である．
従って$\Db(\cC)\to\Db(\AR)$は完全忠実でない．
また$p=n$なら$P$は自然な導来切断で消えるが，genuine avatarでは
$\alpha+\beta$が零指数根となるためこの同時共鳴自体が排除される．

\subsection{反標準側の二つのcharacter}
\begin{example}[実在する指数$(2,-1)$]
滑らかな平面三次$E$の変曲点を原点に取り，非torsion $q$について
\[
(p_1,\ldots,p_6)=(q,3q,2q,5q,7q,11q)
\]
とする．三点共線なら対応する正整数の和の$q$倍が零となり，六点conicも同様なので，
六点は一般位置である．$\alpha=e_1-e_2$，$\beta=e_2-e_3$の指数は$2,-1$．
Rainsのsimple-root effectivityと既約sheaf存在結果
\cite[\S4]{Filtered}\cite[\S11, p.~141]{Rains}の適用条件を
補題\ref{lem:Rains-simple-existence}で確認する．このchartでは上記の$S_+,S_-$が同じheartに実在する．
数値計算だけをsheaf存在の代わりには用いない．
\end{example}

\begin{lemma}[境界外でのExtの比較]\label{lem:localization}
有限次元$A=R/(g-1)$加群$U,V$を境界外の$\qgr R$の対象とみなすと，
全$j\ge0$で$\Ext_A^j(U,V)\simeq\Ext_{\qgr R}^j(U,V)$である．
\end{lemma}
\begin{proof}
$R[g^{-1}]$は強次数付き環で，その次数付き加群圏は$\operatorname{Mod}A$と同値である．
exactな中心局所化の右随伴はrestrictionで，これもexactかつ完全忠実である．
局所化加群は内部次数torsionに導来的に直交するため，右随伴は$\QGr R$にも降りる．
よって導来随伴も完全忠実となる．有限次元加群の非負次数格子は$k[g]$上有限自由で，
有限生成$R$加群を与える．locally noetherian圏でのcoherent Extとの一致により主張を得る．
\end{proof}

\begin{theorem}[三角表現と四次元商]\label{thm:affinequotient}
上の指数$(2,-1)$の例で$V=\Phi(S_+)$，$U=\Phi(S_-)[1]$，$W=\Phi(P)$とする．
これらはそれぞれ次元$2,1,1$の$A$加群であり，
\[
0\longrightarrow W\longrightarrow V\longrightarrow U\longrightarrow0
\]
は非分裂で，$W\not\simeq U$である．$A$の$V$への作用像は全三角行列代数となる．
逆向き非分裂拡大$0\to U\to V'\to W\to0$も存在し，
\[
A/\Ann_A(V\oplus V')\simeq\Pi(A_2).
\]
新しい二単純の根類は$\alpha+\beta,-\beta$であり，指数は$1,1$である．
\end{theorem}
\begin{proof}
最大rankから短完全列を得る．定理\ref{thm:derived}，補題\ref{lem:localization}と
$\Ext_X^1(P,S_-)=0$から$\Hom_A(W,U)=0$となり，二つの一次元加群は異なる．
非零包含$S_-\to P$は同値によって非零の$\Ext_A^1(U,W)$へ移るので，拡大は非分裂である．
異なるcharacterを持つ二次元非分裂表現の像には，二つの対角冪等元と非零の非対角cornerがあり，
像は全三角行列代数である．

同様に$\Ext_A^1(W,U)=\Ext_X^2(P,S_-)=\Hom_X(S_-,P)^*=k$から$V'$が存在する．
$V\oplus V'$への像は，共通の二つの対角成分と両向きの非対角成分からなる
四次元代数$C_0\simeq\Pi(A_2)$の部分代数$C$である．characterの分離から
$C$の半単純像は$k\times k$である．$J=\operatorname{rad}C_0$は$J^2=0$なので，
冪等元を持ち上げ，$1+J$による共役で標準の対角冪等元に合わせる．
両拡大の非分裂性により二つの一次元cornerがともに像へ入り，$C=C_0$となる．
最後に$[P]=[S_+]+[S_-]$，$[S_-[1]]=-[S_-]$を使う．
\end{proof}

反射$s_\beta(v)=v+(v\beta)\beta$は$(\alpha,\beta)$を$(\alpha+\beta,-\beta)$へ送る．
従ってこの局所模型では，正負の共鳴から反標準heartの正指数二単純への組替えがWeyl反射に対応する．
ここまでの主張は二根部分系に限る．Rainsの単純根存在結果を適用した$E_6$への拡張は，
後の系\ref{cor:E6spectrum}で別に与える．
商$A\to\Pi(A_2)$もhomological epimorphismではない．商の単純加群の高次Extは二周期的だが，
ambientのExtは次数2を超えて消えるからである．

\begin{remark}[同型型と生成元の係数]
この定理は実際の有限次元商の同型型を与える．しかし，楕円データからあらかじめ固定した
$A$の生成元に対するcharacterの値や拡大cocycleの係数を数値式で与えてはいない．
その明示計算と全$E_6$のアフィン作用行列は別の課題として残る．
単純加群の根類と次元の決定とは区別する．
\end{remark}

\section{反標準Euler標数によるheartの選択}\label{sec:rootchambers}
ADE型のspherical対象とWeyl反射，および単純tiltによる隣接heartの記述は既知の機構である．
Kleinian特異点についてはBridgeland \cite[\S2, Lemma 3.5]{BridgelandKleinian}を参照されたい．
本節の目的は，その機構を反標準比較へ適用し，\emph{指定構造対象のEuler標数}が
実際の標的heartをどのように選ぶかを明らかにすることである．
内部の$\Db(\cZ)$をambient導来部分圏と同一視する必要はなく，zigzagの形式性も仮定しない．
一般のlength heartの単純tiltが再びlengthとなる条件は\cite[Theorem 2.7]{BCPW}で扱われている．
ここでは2-spherical・単純辺という仮定から必要な有限長性を直接証明する．
幾何学的適用には第\ref{sec:A2}節と同じgenuine楕円gluing・非torsionの仮定を用いる．

\subsection{反標準側の有限次元heart}
$\cD_{\cZ}\subset\Db(\AR)$を，旧heartでの全cohomologyが$\cZ$に属する対象の充満三角部分圏とする．
定理\ref{thm:derived}の自然な$\Phi$を用い，$A=R/(g-1)$とする．
以下の$\operatorname{mod}_{\rm fd}A$は有限次元加群の\emph{abel圏}であり，
その内在的有界導来圏を$\cD_{\cZ}$と同定しているのではない．

\begin{lemma}\label{lem:targetfiniteheart}
$\Phi$は$\cD_{\cZ}$を，境界非交叉かつcohomologyが有限次元$A$加群に対応する
標的の充満三角部分圏へ同値に写す．従って$\cD_{\cZ}$には二つの有界heart
\[
 \cZ,\qquad \cK=\Phi^{-1}(\operatorname{mod}_{\rm fd}A)
\]
があり，$\cZ\subset D^{[0,1]}_{\cK}$である．
この部分圏でSerre関手は自然に$[2]$となり，
$m(X)=\chi(e,X)$は$0\ne X\in\cK$で正整数となる．
\end{lemma}
\begin{proof}
境界triangleの自然な射$\theta\to\mathrm{id}$は$\cD_{\cZ}$上で同型となる．
$\Phi$はこの射を中心元$g$に送るので，像のcohomologyは$g$可逆である．
各$M\in\cZ$について$\Ext^j(e,s^nM)$は$n$に関して周期的で有限次元だから，
像の二つのcohomologyは有限次元$A$加群となる．有界拡大により全$\cD_{\cZ}$で同様である．

逆に有限次元$A$加群は$k[g]$上有限自由な非負次数格子から$\qgr R$の対象を与える．
その逆像$X$には$\theta X\simeq X$が自然な境界射から従う．$\theta$は旧heartにexactなので
全$H^i(X)$も境界非交叉である．$\theta$の数値作用で$c_1$はrank倍の$K_X$だけ変わるから，
$K_X^2>0$の下でそのrankは零．Rainsの次元関数のrank-zero部分は次元高々1なので
$H^i(X)\in\cZ$となる．これは有限次元cohomologyを持つ標的の有界対象にも拡張される．
従って両t構造はこの充満部分圏へ制限される．
$\theta\simeq\mathrm{id}$からSerre関手は$[2]$．
中心局所化の導来随伴を$\pi R$と有限次元$U$に適用すると
$\RHom(\pi R,U)=\RHom_A(A,U)=U$であり，$m(\Phi^{-1}U)=\dim_k U>0$を得る．
\end{proof}

\subsection{全heartのずれは境界非交叉部分で検出される}
\begin{lemma}[中心元の全射性と有限次元局所化]\label{lem:g-surjection}
$U\in\qgr R$について，境界への通常の制限$i_R^*U$が零なら，
$g:U(-1)\to U$は同型であり，$U$は有限次元$A=R/(g-1)$加群から来る．
\end{lemma}
\begin{proof}
有限生成の次数付き代表$M$を取る．$i_R^*U=0$なら$M/gM$は有限次元なので，
$n\gg0$で$g:M_{n-1}\to M_n$は全射である．有限非負整数$\dim M_n$はやがて単調非増加となり，
従って定数になる．十分高い次数では$g$は同型となり，その核も有限次元である．
ゆえに$g$は$\qgr R$で同型で，$\varinjlim(M_n,g)$は有限次元．
この極限が$(M[g^{-1}])_0$であり，強次数付き環による局所化が所要の$A$加群を与える．
\end{proof}

\begin{theorem}[全heartの差を検出する負指数単純根]\label{thm:globaldefect}
定理\ref{thm:derived}，\ref{thm:globaltilt}のgenuine楕円gluingで$q$は非torsionとする．
全Rains heartのtorsionfree部分$\cF=\ker H^0\Phi$は$\cZ$に含まれる．
従って次の三条件は同値である：
\[
 \Phi\text{ はt-exact}
 \quad\Longleftrightarrow\quad(e,s)\text{ はAZ-ample}
 \quad\Longleftrightarrow\quad
 \chi(e,S)>0\text{ for every simple }S\in\cZ.
\]
特にこの範囲では，t-exactnessの失敗は実在する負指数の境界非交叉単純根によって必ず検出される．
\end{theorem}
\begin{proof}
定理\ref{thm:derived}の構成は$\Phi i_{q*}\simeq i_{R*}$を与えるので，左随伴の一意性から
$Li_R^*\Phi\simeq Li_q^*$となる．$M\in\cF$を$\Phi M=U[-1]$と書く．
旧heartでは境界の押出しがexactかつfaithfulであり\cite[\S5冒頭]{Filtered}，
境界triangleと$\theta$のexactnessから$Li_q^*$は振幅$[-1,0]$を持つ．
標的の$Li_R^*$も中心正則元の長さ1の解消から同じ振幅を持つため，
$Li_q^*M\simeq Li_R^*U[-1]$は次数0だけにcohomologyを持つ．
特に$H^0(Li_R^*U)=0$である．補題\ref{lem:g-surjection}から$Li_R^*U=0$となり，
$U$は有限次元$A$加群となる．補題\ref{lem:targetfiniteheart}より$M\in\cZ$である．

$\cZ$の単純対象はtorsion分解によって$\cT$か$\cF$の一方に入る．
そのEuler値はそれぞれ正，負であって零ではない．
$\cF\ne0$なら，その非零対象の$\cZ$での単純部分対象も$\cF$に入り，負指数を持つ．
逆に負指数単純があれば$\cF\ne0$である．t-exactnessと$\cF=0$の同値を用いる．
$t$-exactならpointedな同値が標準ample pairを運ぶのでAZ-ampleである．
逆に非零$M\in\cF\subset\cZ$では$\Hom(e,s^nM)=0$が全$n$で成り立ち，AZ生成性に反する．
\end{proof}
\begin{remark}
ここで得た逆向きの含意は，指定genuine gluing，非torsion，既知の導来比較と振幅を使う．
数値的共鳴の有無だけを判定しているのではなく，\emph{実在する向き付き単純根}を全て調べる条件である．
対象の構成やこれらの仮定を外した一般の非可換曲面への主張ではない．
\end{remark}

\subsection{全ての部分対象・商によるtorsion pairの認識}
\begin{proposition}[Euler標数によるtorsion pairの完全な記述]\label{prop:numericaltorsion}
$\cZ$上の反標準torsion pairは
\[
\begin{split}
 \cT\cap\cZ&=\{M:\ m(Q)>0\text{ for every nonzero quotient }Q\text{ of }M\},\\
 \cF\cap\cZ&=\{M:\ m(N)<0\text{ for every nonzero subobject }N\text{ of }M\}
\end{split}
\]
である．零対象に対する条件は空条件とする．
\end{proposition}
\begin{proof}
非零$T\in\cT\cap\cZ$では$\Phi T$が非零有限次元加群なので$m(T)>0$．
非零$F\in\cF\cap\cZ$では$\Phi F=U[-1]$で$m(F)=-\dim U<0$．
$\cT$は商に，$\cF$は部分対象に閉じるため，それぞれ表示した条件を満たす．
逆に$0\to T\to M\to F\to0$をtorsion分解とすれば，商の条件は$F\ne0$を排除し，
部分対象の条件は$T\ne0$を排除する．
\end{proof}
$m(M)$の符号だけでは不十分である．混合$A_2$の$P,Q$が同じEuler標数を持つ例は，
上の全商・全部分対象の条件を全対象の符号へ置き換えられないことを示す．

\subsection{有限ADE型における単純tiltの有限手順}
まず一般のExt有限な三角圏$\cD_0$でSerre関手が$[2]$である場合を考える．
その有界有限長heart $\cH$の単純対象を$S_1,\ldots,S_r$とし，
全て2-sphericalで，異なる単純の$\Ext^1$が有限ADE図式の隣接数に等しいと仮定する．
$\alpha_i=[S_i]$，$(\alpha_i,\alpha_j)=-\chi(S_i,S_j)$を負定値根形式とする．

\begin{lemma}[一つの単純対象でのtilt]\label{lem:simpletilt}
頂点$j$について
$\cF_j=\add(S_j)$，$\cT_j=\{M:\Hom(M,S_j)=0\}$とすると，$(\cT_j,\cF_j)$はtorsion pairである．
$\cH'=\langle S_j[1],\cT_j\rangle$は有限長heartで，その単純対象は
\[
 S_j[1],\qquad
 P_i\ (i\text{ が }j\text{ に隣接}),\qquad
 S_i\ (i\ne j\text{ が }j\text{ に非隣接})
\]
である．ここで$P_i$は非分裂拡大$0\to S_j\to P_i\to S_i\to0$．
新しい単純根は$s_j(\alpha_i)$であり，$s_j(v)=v+(v,\alpha_j)\alpha_j$とする．
新しい単純も2-sphericalで，同じADE図式を持つ．
\end{lemma}
\begin{proof}
$\Ext^1(S_j,S_j)=0$なので，最大の$S_j$-半単純商への評価射
$M\twoheadrightarrow S_j^{\oplus\dim\Hom(M,S_j)}$の核は$\cT_j$に入る．
これがtorsion分解を与える．新heartの候補は表示した通りで，隣接$i$については
新heartで$0\to P_i\to S_i\to S_j[1]\to0$となる．

有限長性を確かめる．$M\in\cT_j$の旧heartでの単純商を$S_i$とすると$i\ne j$である．
核$K$のtorsion部分を$T_j(K)$とし，$E=M/T_j(K)$と置く．
$0\to S_j^{\oplus b}\to E\to S_i\to0$であり，$\Hom(E,S_j)=0$から
$k^b\hookrightarrow\Ext^1(S_i,S_j)$を得る．従って$b=0$，または$i$が隣接して$b=1$の非分裂拡大だけである．
旧lengthによる帰納法と上の新heartの短完全列により，$M$は表示した候補で有限filtrationを持つ．
全新heartもそのようなfiltrationを持つ．

候補間のHomは異なる二対象で零，同じ対象で$k$である．
例えば$\Hom(P_i,S_j)=0$は非零拡大類によるconnecting mapから従い，
$\Hom(P_i,S_j[1])=0$はSerre双対性とEuler形式から得る．
異なる$P_i,P_k$間では，$P_k$の部分対象$S_j$と商$S_k$へのHomを用いれば零となる．
従って候補は本当に全単純対象である．その根類は表示した反射で，反射はEuler形式を保つ．
Serre双対性とHomの計算から各自己Ext$^1$は零，cross Ext$^1$は元の隣接数となる．
\end{proof}

\begin{theorem}[反標準heartのEuler chamber認識]\label{thm:Eulerchamber}
$\cH\subset\cD_0$を上の有限ADE型heartとし，もう一つの有界heart$\cK$について
$\cH\subset D^{[0,1]}_{\cK}$とする．加法的関数$m:K_0(\cD_0)\to\mathbb Z$が
全ての非零$U\in\cK$で$m(U)>0$を満たすと仮定する．
このとき次の有限手順は$\cK$を回復する．
\begin{enumerate}
\item 現在の単純のうち$m(S_j)<0$のものを一つ選ぶ．
\item 補題\ref{lem:simpletilt}の$\langle S_j[1],\cT_j\rangle$へ移る．
\item 負の単純がなくなるまで繰り返す．
\end{enumerate}
手順は初期正根系$\mathcal R_+$について
\[
 N(m)=\#\{\alpha\in\mathcal R_+:m(\alpha)<0\}
\]
回で終了する．終端heartは選択順序に依存せず$\cK$であり，その単純根系は
正根系$\{\alpha:m(\alpha)>0\}$の単純根系に等しい．
特に$m$は全ての根で非零であることが従う．
\end{theorem}
\begin{proof}
現在のheartの単純$S$には，振幅1の比較からtorsion分解がある．単純性により
$S\in\cK$または$S\in\cK[-1]$のいずれかであり，その$m$はそれぞれ正または負である．
従って$m(S)=0$は起こらない．負の$S_j$を選ぶと$S_j[1]\in\cK$であるから，
新heartも$D^{[0,1]}_{\cK}$に入る．補題により有限長性とADE条件も保たれる．

現在の正根系を単純反射$s_j$で反射すると，正根の集合は$\alpha_j$だけを
$-\alpha_j$に入れ替えたものとなる．その$m$は負から正へ変わるので，負の正根の数は正確に1減る．
この数は非負整数であるため手順は停止する．全単純の$m$が正となったとき，
それらは$\cK$に入り，新heartの有限長性から新heart全体が$\cK$に含まれる．
同じ三角圏の有界heartの包含は一致を意味するので，終端は$\cK$である．
終端の各正根は単純根の非負整数結合で$m>0$となり，全根での非零性も従う．
最後に，正根系を指定するregular線形汎関数のchamberは一意なので単純根系も一意である．
対象としての終端heartの独立性は，根類の一致だけでなく，既に示した$\cK$との一致による．
\end{proof}
\begin{remark}
これは一般のWeyl反射や単純tiltを新しく導入する定理ではない．
本稿では特に$\cH=\cZ$，$m=\chi(e,-)$，$\cK=\Phi^{-1}(\operatorname{mod}_{\rm fd}A)$に適用し，
\emph{反標準比較が指定する}heartと，そこでの有限次元単純加群の次元を決める．
得られた根類$\beta_i$に対応する単純加群の次元は$m(\beta_i)$である．
異なる経路が同じheartに到達することと，それらの合成自己同値が自然に同型であることは別であり，後者は主張しない．
またこれは固定圏におけるEuler汎関数によるheartの選択であり，
楕円パラメータの族での全ての壁越えや安定性空間全体の記述ではない．
\end{remark}

\subsection{局所的な反射手順の全heartへの延長}
\begin{lemma}[境界非交叉単純の有限ADE性]\label{lem:finiteADEautomatic}
定理\ref{thm:globaldefect}の設定では，$\cZ$の単純対象は有限個で，
その根類は線形独立な有限ADE単純根系をなす．空集合も許す．
\end{lemma}
\begin{proof}
単純$S$に対して$\theta S=S$であり，$\End(S)=\Ext^2(S,S)=k$．
$c_1(S)=\alpha\ne0$，$H\alpha=0$と負定値性，Euler形式の偶奇から
$0<-\alpha^2=2-\dim\Ext^1(S,S)\le2$，従って$\alpha^2=-2$である．
同根の異なる単純はHomとExt$^2$が零なのにEuler値が2となって矛盾する．
負定値格子には長さ2の根が有限個しかないので，単純も有限個である．
異なる二根の交点はExt$^1$の次元に等しく非負．
反対根は有効錐のpointednessで排除されるので，交点は0か1である．

根類が有理係数の線形関係を持つとする．正負の係数を分けて$x=y$と書けば，
支持が異なるため$x\cdot y\ge0$，従って$x^2\ge0$となる．
負定値性から$x=0$であるが，非自明な正係数の有効類の和が零であることは
有効錐のpointednessに反する．片側しかない関係も同様に排除される．
従って根類は独立であり，対角2，非対角0または$-1$の正定値Cartan行列を持つ．
有限型の分類からADE型である．
\end{proof}

\begin{theorem}[全Rains heartの有限単純tilt分解]\label{thm:globalADEtilts}
定理\ref{thm:globaldefect}の範囲で，定理\ref{thm:Eulerchamber}の負指数単純tiltは
全Rains heart上に延長できる．初期$\cZ$の正根系を$\mathcal R_+$とすると，
$\#\{\alpha\in\mathcal R_+:m(\alpha)<0\}$回で全反標準heartへ到達する．
単純対象がない場合は初めから両heartは一致する．
\end{theorem}
\begin{proof}
$\cZ$がSerre部分圏なので，その単純対象は全heartでも単純である．
$\Hom$の有限性と$\Ext^1(S_j,S_j)=0$から，全対象に対する最大$S_j$-半単純商を用いて
$(\ker\Hom(-,S_j),\add S_j)$というtorsion pairを全heart上で構成できる．
負指数なら$S_j[1]$は標的heartに属すから，tilt後も振幅$[0,1]$を保つ．

現在の境界非交叉Serre部分圏を$\cZ'$とする．新heartの短完全列の旧cohomologyを取ると，
次数$-1$の項は全て$\add S_j\subset\cZ'$に入る．従って旧$\cZ'$のSerre性から，
$\cD_{\cZ}$と新heartとの交わりもSerre部分圏である．その制限は補題\ref{lem:simpletilt}のheartである．
新しいtorsionfree部分の対象$X$を，旧cohomologyのtriangle
$H^{-1}(X)[1]\to X\to H^0(X)$で記述する．
$H^{-1}(X)[1]$は標的heartに属し，$H^0(X)$の標的cohomologyは負次数にないため，
$H^{-1}(X)[1]\hookrightarrow H^0_{\rm target}(X)=0$となる．
よって$X$は旧heartにも属し，旧torsionfree部分の対象である．
従って各段階の残るtorsionfree部分は$\cD_{\cZ}$の中にある．

有限ADE性は補題\ref{lem:finiteADEautomatic}で得られる．局所的な手順が標的の有限次元heartに
到達すると，残るtorsionfree部分は零になる．命題\ref{prop:hrs}により全heartも一致する．
回数は定理\ref{thm:Eulerchamber}から従う．
\end{proof}
この有限分解は一回のHRS tiltという全体記述を，有限個のspherical単純による操作へ細分する．
中間heartを別に指定したblowdown表示のRains heartと同一視する必要はない．

\subsection{最大共鳴の単純対象}

\begin{proposition}[全単純根が実在する場合の網羅性]\label{prop:maximalrootexhaustion}
$\cZ$の根格子が有限ADE型で，選んだ全単純根$\alpha_i$に対する
既約境界非交叉sheafが実在すると仮定する．このときそれらが$\cZ$の全単純対象であり，
同根の単純は同型を除き一意である．
\end{proposition}
\begin{proof}
任意の単純$M\in\cZ$は根類を持つ．実際，$D=c_1(M)\ne0$は$H^\perp$にあり，
負定値性，Riemann--Rochの偶奇，Serre双対性から
$0<-D^2=2-\dim\Ext^1(M,M)\le2$となるので$D^2=-2$である．
負根が有効なら，その反数が実在単純根の非負和として有効になり，
有効錐のpointedness \cite[Proposition 7.19]{Rains}に反する．
非単純正根$\gamma=\sum c_i\alpha_i$では
$-2=\gamma^2=\sum c_i(\gamma,\alpha_i)$だから，ある$i$で$(\gamma,\alpha_i)<0$．
しかし異なる単純間のExt$^1$の次元はこの交点に等しいため不可能である．
同根の異なる単純も，HomとExt$^2$が零なのにEuler値が2となって矛盾する．
\end{proof}

\subsection{単純根の実在入力と整数倍点族の完全な次元表}
\begin{lemma}[Rainsの単純根存在結果の適用]\label{lem:Rains-simple-existence}
滑らかな楕円境界を持つ指定Rains blowdown表示で，$\alpha$をその向き付き単純根とする．
$q$が非torsionで，ある整数$m$について$\rho(\alpha)q^m=\OO_E$なら，
境界非交叉の既約純一次元sheaf $S_\alpha$で
$c_1(S_\alpha)=\alpha$，$\chi(e,S_\alpha)=m$となるものが存在する．
本稿のnef・正次数の範囲では，これは単純であり同型を除き一意である．
\end{lemma}
\begin{proof}
存在はRainsの既知結果を適用する．\cite[\S4, pp.~34--35]{Filtered}は，
指定blowdownの単純根が$q$の冪へのrestrictionを持つとき有効となることと，
対応する$(-2)$類のpure sheafに対するdeterminant条件を記述している．
\cite[\S11, p.~141, Proposition 11.23の直前]{Rains}では，
\emph{有効単純根そのもの}の類に対し，そのdeterminant条件を満たすEuler標数で
既約sheafが存在し，境界と交わらないことが明記されている．
Proposition 11.23のnef・fundamental chamberの主張を単純根に誤用しているのではない．
本稿の規約でrank-zero類のrestriction determinantは$\rho(\alpha)q^m$である．
従って入力の仮定を満たす．単純性は補題\ref{lem:offboundary}，一意性は
同根の異なる単純間でEuler値2となる矛盾による．
\end{proof}
\begin{remark}
これは任意の数値的共鳴根に無条件にsheafを割り当てる主張ではない．
\emph{固定したRains heartのblowdown単純根}であることを追加して既知の存在結果を使っている．
以下の整数倍点族では，各単純根についてこの条件を個別に確認する．
\end{remark}

\begin{theorem}[整数倍点族：整列した差による全単純の次元]\label{thm:sorted-spectrum}
$E\subset\mathbb P^2$のflexを原点$o$とし，$q\in E$を非torsionとする．
$3\le r\le8$，相異なる正整数$N_1,\ldots,N_r$に対して$p_i=N_iq$と置く．
$S=\operatorname{Bl}_{p_1,\ldots,p_r}\mathbb P^2$と，この順序のblowdownによるRains heartを取る．
$q$のPicard類は$\OO_E(q-o)$であり，共鳴の符号は\eqref{eq:resonance}に従う．
このとき$S$はgenuine del Pezzoであり，境界非交叉heartの全単純根は
\[
 \alpha_0=h-e_1-e_2-e_3,\qquad \alpha_i=e_i-e_{i+1}\quad(1\le i<r)
\]
である．対応するsheafは全て実在し，その指数は
$(N_1+N_2+N_3,N_2-N_1,\ldots,N_r-N_{r-1})$である．

$N_{\pi(1)}<\cdots<N_{\pi(r)}$となる並べ替えを$\pi$とする．
$A=R/(g-1)$の有限次元単純加群は同型を除き正確に$r$個で，その根類と次元は
\begin{align*}
 \beta_0&=h-e_{\pi(1)}-e_{\pi(2)}-e_{\pi(3)},
 &\dim V_0&=N_{\pi(1)}+N_{\pi(2)}+N_{\pi(3)},\\
 \beta_i&=e_{\pi(i)}-e_{\pi(i+1)},
 &\dim V_i&=N_{\pi(i+1)}-N_{\pi(i)}\quad(1\le i<r).
\end{align*}
全Rains heartから反標準heartへの負単純tiltの回数は
\[
 \operatorname{inv}(N)=\#\{(i,j):i<j,\ N_i>N_j\}
\]
である．特に両heartが一致することと$N_1<\cdots<N_r$は同値である．
\end{theorem}
\begin{proof}
$q$が非torsionなので点は相異なる．blow-up後の滑らかな反標準曲線の類$H$はnefで，
$H^2=9-r>0$である．$H^\perp$の根を$ah-\sum b_ie_i$と書くと
$\sum b_i=3a$，$\sum b_i^2=a^2+2$．$a\ge0$の根は差根，三点の和を引く直線類，
六点の和を引くconic類，および$r=8$で一つの係数が2・残りが1の三次曲線類だけである．
実際Cauchy--Schwarzで$(9-r)a^2\le2r$，かつ
$\sum b_i(b_i-1)=a^2-3a+2$を使えば$a\le3$とこの列挙を得る．
可能性$a=4,r=8$では右辺6を一つの3または三つの2で分配しても$\sum b_i=12$に届かず排除される．

restrictionは$\rho(ah-\sum b_ie_i)=q^{-\sum b_iN_i}$である．
差根の指数は相異なる二整数の差，$a>0$の根の指数は正整数の正係数和なので零にならない．
もし$H$がampleでなければ，nef・bigとNakai--Moishezonから$HC=0$の既約曲線がある．
Hodge指数定理とadjunctionより$C^2=-2$で，境界と交わらずrestrictionが自明となり矛盾する．
よって$S$はgenuineである．全ての表示した単純根は$q$の整数冪にrestrictionするため，
補題\ref{lem:Rains-simple-existence}を\emph{同じblowdown表示}で各根に適用する．
命題\ref{prop:maximalrootexhaustion}が全単純の網羅性を与える．

正根のうち負指数を持つものは差根$e_i-e_j$ $(i<j,N_i>N_j)$だけである．
負の隣接差を選ぶ単純反射は隣接二整数の交換であり，$\alpha_0$の指数は常に
現在の先頭三整数の和だから正である．従って隣接逆転を解消する整列が
定理\ref{thm:Eulerchamber}の手順となり，終端の単純根と正指数は表示したものになる．
補題\ref{lem:targetfiniteheart}がこれらを全有限次元単純加群に対応させる．
定理\ref{thm:globalADEtilts}から同じ回数で全heartも一致する．
初期が整列済みのときだけ負の単純がないので，最後の同値も従う．
\end{proof}

\begin{corollary}[同じpointed導来圏内の異なるRains hearts]\label{cor:permuted-hearts}
上の無順序点集合を固定すると，その$r!$通りの順序から得るRains heartsは，
同じ$\cD_q$と同じ$e$の中で相異なる．反標準環$R(e,s)$は共通であり，
この$r!$個のうち反標準heartと一致するのは整数を昇順に並べたものだけである．
\end{corollary}
\begin{proof}
可換blow-up $S$は点の順序によらず，gluing\eqref{eq:glue}は$(S,E,q)$と$e$から定まる．
従って導来圏と構造対象，Serre関手，反標準環は共通に固定できる．
一方，異なる二順序には前後関係が逆になる点の対がある．一方のheartでは$e_i-e_j$が
隣接単純根の正和として有効で，他方ではその反数が有効である．
もし両heartが同じなら有効錐のpointednessに反する．最後は定理\ref{thm:sorted-spectrum}による．
\end{proof}

\begin{corollary}[六点の具体例]\label{cor:E6spectrum}
$(N_1,\ldots,N_6)=(1,3,2,5,7,11)$では，全六単純根sheafが実在し，
一回の$\alpha_2$-tiltで全反標準heartへ移る．全有限次元単純の根類と次元は
\[
\begin{array}{c|rrrrrr}
\text{根類}&\alpha_0&\alpha_1+\alpha_2&-\alpha_2&\alpha_3+\alpha_2&\alpha_4&\alpha_5\\\hline
\text{次元}&6&1&1&2&2&4
\end{array}
\]
であり，一次元characterは正確に二つである．同じ無順序六点集合は同じpointed導来圏内に
$6!=720$個の相異なるRains heartsを与える．
\end{corollary}
\begin{proof}
整列は$(1,2,3,5,7,11)$で，逆転は一つだけである．前二結果を適用する．
\end{proof}
\begin{remark}[共鳴とheartの不一致は同義ではない]
この族では全根が数値的に共鳴しているが，整列した順序のRains heartは反標準heartと一致し，
$(e,s)$はAZ-ampleである．非整列順序では少なくとも一つの負指数単純があり，AZ-ampleでない．
従って最大の数値的共鳴とheartの一致は両立する．この主張はheartの\emph{同じ導来圏内での位置}についてであり，
構造対象を忘れた抽象的abel同値の全分類ではない．
単純加群の根類・次元は決定したが，元の楕円生成元の作用行列や全アフィン環の表示は未計算である．
既知のRains存在結果とWeyl反射を適用した帰結として扱い，有限次元表現論一般の新発見とは主張しない．
\end{remark}


\section{四者を同じ図式に置くための条件}\label{sec:alignment}
比較は二圏ずつの同値の存在だけでは完了しない．同じ$e$と$s$を使って反標準環を作り，
比較の合成がその構造を保つことを確認する必要がある．
\begin{proposition}[pointedな比較の接続]\label{prop:alignment}
次のデータが与えられたとする．
同じblow-up表示によるpointedな$J:\Db(\AV)\simeq\Db(\AR)$，
定理\ref{thm:derived}の$\Phi$，およびpointedな
$E:\Db(\AAS)\simeq\Db(\AR)$がある．さらにAS側で系\ref{cor:ASac}の条件と
命題\ref{prop:unique}の高次消滅が成立するとする．
このとき$\Phi E$は標準AS heartから反標準heartへのt-exact同値であり，
$\Phi J$はVdB heartに対して$J^{-1}(\cT),J^{-1}(\cF)$によるHRS比較となる．
従って$\cF=0$なら四者はこれらの指定関手の下で一致する．
\end{proposition}
\begin{proof}
corner同値からAS側の指定反標準orbitは豊富である．pointed性により環とorbitを同定でき，
両heartは命題\ref{prop:unique}の認識式を満たすので$\Phi E$はt-exact．
$J$のt-exactnessはtorsion pairを$\AV$へ運び，定理\ref{thm:globaltilt}の式を保つ．
$\cF=0$のときの結論は命題\ref{prop:hrs}から従う．
\end{proof}

実際のrestriction条件\eqref{eq:two-nontrivial}の下では，定理\ref{thm:AOU-pointed}により
section model--AS比較のpointed性が得られる．
VdBの最後のraw blow-upと同じRains表示については，系\ref{cor:fourfunctors}で
指定関手を実際に接続した．異なるavatarとの比較にはなおデータの同定を要する．
確認済みの範囲と残る同定は，次の表のとおりである．
\begin{center}
\small
\begin{tabular}{p{.18\textwidth}p{.36\textwidth}p{.35\textwidth}}
\toprule
比較 & 結論と適用範囲 & 自動的には従わないこと\\
\midrule
VdB--Rains & 同じblow-up・gluing・blow-down表示でabel同値 & 独立に選んだavatarのデータ対応\\
Rains--反標準 & genuine smooth楕円gluingで導来同値，一回のHRS tilt & 任意nc曲面・特異境界への拡張\\
VdB--section model & VdBの一般位置条件下で振幅1の導来同値 & raw blow-upとのheart一致\\
AS--反標準 & 規格化共通領域でnoetherian性と標準heartの一致を導出 & 任意potential・全AS locusへの拡張\\
section model--AS & 実際のrestriction条件の下でpointed導来同値，反標準環も一致 & 全域での標準heartの一致\\
四者の接続 & 同じ最終VdB表示では自然関手を同定，$\Delta\ne0$で単一tilt & 独立avatarの同定，相対降下\\
\bottomrule
\end{tabular}
\end{center}

\subsection{可換基準点と非可換共鳴}
滑らかな可換三次曲面$S$では，実際のblow-up heartとRainsの可換heartは$\coh S$である．
$-K_S$がampleなので反標準環の$\qgr$も$\coh S$を回復する．
標準的なline bundle helixの各周期部分列$\OO(D_i-nH)$もSerreの生成・消滅を満たし，
ASの$\mathbb Z$-代数から同じheartを再構成する．この場合には四者の一致を実際に確認できる．
しかし，一つの可換点での一致から，非可換パラメータ空間での一致locusの開性を推論しない．
命題\ref{prop:density}は，楕円chartでは負指数共鳴が可算個の軌跡の合併として稠密となることを示す．
「稠密な共鳴」と「一つの閉じた判別因子」，あるいは「導来特異化」は異なる概念である．

\subsection{族に移すときの追加課題}
本稿の主要な比較は幾何学的点での結果である．基礎$T$上の族へ移すには，
少なくともgluing bimoduleとkernelの相対perfect性，orbitのHomの基底変更，
有限性と消滅次数の一様性，t構造の基底変更適合性を確認する必要がある．
構造対象の同型と反Serre作用への適合を降下データとして保つことも必要である．
構造対象の同定については命題\ref{prop:AOU-relative-object}が相対Homのperfectnessと
基底変更の下で評価triangleから自然な同型を与える．その前提となる相対圏の構成と
全heartの降下を，この対象単独の同定から推論しない．

有限cornerの範囲だけなら次の初等的な相対判定が使える．
\begin{proposition}\label{prop:relativecorner}
有限型基礎$T$上に次数1生成のgraded代数族$C$があり，各$C_N$が有限表示で
fiberの形成と両立するとする．このときfiberが有限次元となるlocusは
$\bigcup_{N\ge1}\{t:(C_N)_t=0\}$であり開である．
$T$がnoetherianなら，この開locus上である共通の$N$を選べる．
\end{proposition}
\begin{proof}
各集合はcoherent module $C_N$のsupportの補集合であり開である．
補題\ref{lem:cutoff}によってこれらは$N$とともに増加し，合併が求めるlocusとなる．
noetherian空間の開集合は準コンパクトなので有限部分被覆を選び，最大の$N$を取る．
\end{proof}
三次曲面型では第\ref{sec:explicitcorner}節の十一条件がこの開集合の明示式を与える．
この命題は有限cornerのfiber条件を述べるだけであり，heartの相対同値・flatness・stack降下を
証明するものではない．本稿のt構造比較と，内在的非可換del Pezzoモジュライの
全成分分類は分離して扱う．

\appendix
\section{三次曲面の数値データと再現手順}\label{app:numerics}
可換三次曲面の標準格子を
$\Pic(S)=\mathbb Zl\oplus\bigoplus_{i=1}^6\mathbb Ze_i$，
$l^2=1$，$e_i^2=-1$，$H=3l-\sum e_i$とする．
AOU \cite[(4.1)]{AOU}の九線束のChern類は
\[
0,\quad -2l+\sum e_i,\quad -l+e_1+e_2+e_3,
\quad e_1,e_2,e_3,\quad l-e_4,l-e_5,l-e_6.
\]
各三つが一blockをなし，$H$次数は$0,1,2$である．
$\chi(\OO(D),\OO(E))=1+\tfrac12(E-D)(E-D+H)$からEuler行列は
\[
C=\begin{pmatrix}I_3&J_3&2J_3\\0&I_3&J_3\\0&0&I_3\end{pmatrix},
\qquad \det C=1,
\]
ただし$J_3$は全成分1の行列．成分総和は45である．
これは可換点の既知のfull strong性と合わせてEnd代数の次元を確認するが，
Euler行列だけからstrong性を証明してはいない．
$K_0$でのSerre作用は$C^{-1}C^{\mathsf t}$，反Serre作用はその逆であり，
後者は$H$によるtensorと一致する．

根は$e_i-e_j$ $(i\ne j)$，$\pm(l-e_i-e_j-e_k)$ $(i<j<k)$，
$\pm(2l-\sum e_i)$の72個である．各根について$\alpha^2=-2$，$H\alpha=0$．
反標準Euler値は$n=0,\ldots,8$で
\[
1,4,10,19,31,46,64,85,109
\]
となる．これらを添付コードで有理数上の厳密計算として確認する．

混合$A_2$では，四直既約のHom表をquiver射の線形方程式から求め，
Euler形式とSerre双対からambient Ext$^1$を検算する．
指数$p,n=1,\ldots,8$，四直既約の重複度$0,\ldots,4$について
40,000通りのcohomology公式を検査する．$p=n$の検算は仮想的線形代数の確認であって，
genuine avatar上にそのsheafが存在することの証拠ではない．

構造対象のmutationについては，AOUの操作の合成順序を実装し，
標識対象の位置が$9,9,9,6,6,3,1$と移ることと，最後の二箇所だけに
Hom消滅の入力が必要であることを検査する．
実際の$\mathcal U_2,\mathcal U_1$を使った\eqref{eq:two-degree-zero}の次数零性も検査する．
印刷された旧指数の検算は境界規約を保証するものではなく，現在は
第\ref{subsec:boundaryaudit}節の不整合を再現する検査と区別する．
境界検算ではさらに，逆の右mutationの$[-1]$と，末尾三対象への$[1]$を明示して
各段階のrankとdeterminantを計算し，実際の正規化したdecupleの指数と一致することを検査する．
これは非自明な楕円線束のcohomology消滅やmutation対象の同型を計算で証明するものではなく，
証明に使う順序と数値の独立な確認である．

Euler chamberの検算は有限根系$A_2+A_1,A_4,D_5,E_6,E_7,E_8$について行う．
各型で32個のregular汎関数を作り，負の単純を選ぶ三通りの規則を比較する．
停止回数が初期負正根数に一致し，終端根集合が正値根の単純根集合に一致することを整数上で検査する．
これは有限個の根格子計算であって，単純sheafの実在やheartの同定を数値実験に置き換えるものではない．

整数倍点族では$3\le r\le8$の各型で40組の相異なる正整数を選び，
Picard格子での根の明示列挙とWeyl反射による列挙を照合する．
具体的六点の全720順列について，停止回数が逆転数に等しく，全単純の次元が同じ整列公式となることを検査する．
これらは格子の検算であり，Rainsの既知のsheaf存在結果を計算で証明するものではない．

本稿は主TeXファイルと九つのsectionソースからなり，\texttt{latexmk -lualatex}で再生成できる．
日本語LuaLaTeX，LuaTeX-ja，\texttt{tikz-cd}，Python~3を用いる．
リポジトリrootでの主要コマンドは
\begin{verbatim}
python3 computations/verify_cubic_heart_numerics.py
python3 computations/verify_four_hearts.py
python3 computations/verify_four_hearts_corner.py
python3 computations/verify_four_hearts_chambers.py
python3 computations/verify_four_hearts_sorted.py
python3 computations/verify_four_hearts_boundary.py
python3 computations/verify_four_hearts_elliptic.py
latexmk -lualatex papers/nc_del_pezzo_hearts.tex
\end{verbatim}
である．数値検算は，sheafの存在，derived type，導来完全忠実性，
または新規性を計算によって証明するものではない．


